\documentclass[10pt,a4paper]{amsart}
\usepackage[margin=19mm]{geometry}
\usepackage{amsmath,amssymb,mathtools}
\usepackage[scr=rsfs,cal=euler]{mathalfa}
\usepackage{xfrac}
\usepackage[unicode,hidelinks,bookmarks=false]{hyperref}
\newcommand{\ie}{i.e.\ }
\newcommand{\cf}{cf.\ }
\newcommand{\resp}{respectively\ }
\newcommand{\Subsection}{\subsection}
\providecommand{\nobreakdash}{\nobreak\hbox{-}}
\setlength{\emergencystretch}{2em}
\multlinegap0pt
\usepackage{amssymb}
\usepackage{dsfont}
\newtheorem{thm}{Theorem}
\title{Random systems of holomorphic sections of a sequence of line bundles on compact K\"{a}hler manifolds}
\author{Afrim Bojnik, Ozan G\"uny\"uz}
\date{Source: arXiv:2401.08243v7 (CC BY 4.0)}
\begin{document}
\maketitle

\noindent\emph{Source and licence.} Random systems of holomorphic sections of a sequence of line bundles on compact K\"{a}hler manifolds. Version arXiv:2401.08243v7 (CC BY 4.0), distributed by arXiv under the Creative Commons Attribution 4.0 International licence, http://creativecommons.org/licenses/by/4.0/, as recorded in the arXiv licence metadata for this exact version and shown on the abstract page https://arxiv.org/abs/2401.08243v7. Full licence notice: \texttt{licenses/arxiv-2401.08243-CC-BY-4.0.txt}.\par\smallskip

\noindent\textit{Editorial scope.} Counted unit: the article's thm th1 (source0.tex lines 219--225). The article's complete original proof of it is delivered with it (source0.tex lines 664--672, one \texttt{proof} environment of 9 lines, which the article places away from the statement). No mathematics is written, repaired or re-derived: the statement and the proof are the article's own lines, in the article's own notation, and no retained line is substituted or re-flowed.  Retained uncounted as the source-local context the proof needs: lines 503--509.  Nothing was omitted from any retained span, so the package carries no omission record.  No retained line contains a citation, so the wrapper carries no bibliography and no bibliography entry.  No retained line contains a figure environment, an image-inclusion command or drawing code, so no image is delivered and the package declares no asset dependency.  Editorial formatting only, in addition: an AMS \texttt{amsart} preamble that restates the article's own declarations and macros used by the retained lines, namely the environments \texttt{thm} on separate counters, together with the standard packages those lines need.  The retained environments print in this document as Theorem 1 in order of appearance, whereas the article prints the same items with the numbering of its own full document.  The complete native source of the article as distributed in the arXiv e-print for this exact version is bundled unchanged as \texttt{sources/152998/source0.tex}.
\begin{thm} \label{genvar}
Let $L_{p, 1}, \ldots, L_{p, k}, \,\,k\in \{1, \ldots, n\}$, be sequences of $\mathscr{C}^{2}$-Hermitian holomorphic line bundles on a compact K\"{a}hler manifold $(X, \omega)$. Assume that the subspaces $\mathcal{S}_{p, j} \subset H^{0}(X, L_{p, j}),\,j=1,\ldots, k$, equipped with the probability measures $\sigma_{p, j}$ satisfy \textbf{(C1)-(C2)}, and that for every $j=1,\ldots,k$,
the linear system $\mathcal{S}_{p,j}$ is base-point free, i.e. $\operatorname{Bs}(\mathcal{S}_{p,j})=\varnothing$ for large values of $p$. Then for $\Sigma^{k}_{p}=(s_{p}^{1}, \ldots, s_{p}^{k})\in \mathcal{S}_{p, 1} \times \ldots \times \mathcal{S}_{p, k}$ with random sections $s_{p}^{1}, \ldots, s_{p}^{k}$ selected independently with respect to the product probability measure $\sigma_{p, 1} \times \cdots \times \sigma_{p, k}$, we have

\begin{equation} \label{varra}\mathrm{Var} \langle [Z_{\Sigma_{p}^{k}}], \phi \rangle \leq (C_{p})^{2/\alpha}\,(B_{\phi})^{2}(\int_{X}{\omega^{n-k+1} \wedge \sum_{\beta=1}^{k}{\Big[ \prod_{1\leq j \leq k, j \neq \beta}{c_{1}(L_{p, j}, h_{p, j})}\Big]}\,})^{2}\end{equation}

\end{thm}

\begin{thm}\label{th1}
	Let \( (L_{p}, h_{p})_{p \geq 1} \), \( (X, \omega) \), and \( \sigma_{p} \) be as defined above. Assume that they satisfy the conditions \textbf{(A)}, \textbf{(C1)}-\textbf{(C2)},
and the linear system $H^{0}(X, L_p)$ is base-point free, i.e. $\operatorname{Bs}(H^{0}(X, L_p))=\varnothing$ for large values of $p$. Further, let the systems \( \Sigma^{k}_{p}= (s_{p}^{1}, s_{p}^{2}, \ldots, s_{p}^{k}) \) with sections chosen independently with respect to \( \sigma^{k}_p \) be given. Then, there exists \( P \in \mathbb{N} \) such that for all \( p \geq P \) and any \( \phi \in \mathcal{D}^{n-k, n-k}(X) \), the following estimate holds:
	\begin{equation}\label{mainvar}
			\mathrm{Var}\langle [\widehat{Z}_{\Sigma_{p}^{k}}], \phi \rangle  \leq D_{n} \frac{(C_{p})^{2/\alpha}}{A^{2}_{p}}\,(B_{\phi})^{2}(\mathrm{Vol}(X))^{2}
			 \end{equation}where  $B_{\phi}= b \|\phi\|_{\mathscr{C}^{2}}$ is a positive constant depending on the form $\phi$, and $D_{n}$ is a constant that depends only on the dimension of $X$.
\end{thm}

 \begin{proof}[Proof of Theorem \ref{th1}]
 Let us take all holomorphic line bundles as identical and $\mathcal{S}_{p, j}= H^{0}(X, L_{p})$ in Theorem \ref{genvar}. We immediately have

 \begin{equation} \label{spevar}
  \mathrm{Var} \langle [Z_{\Sigma_{p}^{k}}], \phi \rangle \leq (C_{p})^{2/\alpha}\,(B_{\phi})^{2}(n \int_{X}{\omega^{n-k+1} \wedge c_{1}(L_{p}, h_{p})^{k-1}})^{2}.
 \end{equation}By the diophantine approximation \textbf{(A)}, there exists $p_{1} \in \mathbb{N}$ such that
 \begin{equation*}\label{cur1}\frac{A_{p}}{2}\,\, \omega \leq c_{1}(L_{p}, h_{p}) \leq 2A_{p}\,\, \omega \end{equation*} for all $p \geq p_{1}$, which implies that \begin{equation}\label{voll}\int_{X}{\omega^{n-k+1}\wedge c_{1}(L_{p}, h_{p})^{k-1}}\leq (2A_{p})^{k-1} \,\mathrm{Vol}(X) \leq 2^{n-1} A^{k-1}_{p} \mathrm{Vol}(X)\end{equation} for all $p \geq \max\{p_0, p_1\}$. Finally, by using the normalization and writing $D_{n}:= (n \,2^{n-1})^{2}$, Theorem 1.1 follows.

 \end{proof}

\end{document}
